% !PRINTED 444
\weberpage{444}
%<<CONT
one proceeds as if it were a matter of determining the greatest
common divisor of $m$ and $n$.

From 8. and 9. there follows yet one last theorem, which holds when
$n$ and $n'$ are two positive odd numbers coprime to $m$:

\[
10.\quad \left(\frac{m}{n}\right)\left(\frac{m}{n'}\right) = \left(\frac{m}{nn'}\right).
\]

Its correctness follows, when $m$ is odd and positive,
on the ground of the congruence

% [sic] Within §138 the source reuses equation label (7), already printed on p.442 for the alpha formula; preserved diplomatically.
\weq{7}{(n-1)(n'-1) = (nn'-1) - (n-1) - (n'-1) \equiv 0 \ (\mathrm{mod}\ 4),}

if one applies formula 9. and then 8. to both sides of 10.;
and, when $m$ is even or negative, from 4. and 6.

By these theorems one can reduce $\left(\dfrac{m}{n}\right)$ to a product of
values $\left(\dfrac{q}{p}\right)$, in which $p, q$ are primes. The
computation of $\left(\dfrac{m}{n}\right)$, however, is carried out more easily by 9. without this
decomposition.

We now wish to learn another meaning of the symbol
$\left(\dfrac{m}{p}\right)$ in the case where $p$ is an odd prime,
for which \gesperrt{Legendre} originally introduced the symbol.

If we raise to the square all the integers $x$ not divisible by $p$
and seek out the residues of the division by $p$,
we obtain in all only $\dfrac{p-1}{2}$ different
residues; for, first, the squares of congruent numbers give the same
residues, and secondly the squares of opposite
numbers also give the same residues. We therefore certainly obtain all the
residues if we take the $\dfrac{p-1}{2}$ numbers $\nu$

\[
(\nu)\qquad 1,\ 2,\ 3\cdots\frac{p-1}{2}
\]

and raise them to the square. On the other hand, the
squares of these numbers $\nu$ also yield residues that are all different; for,
if $\nu$ and $\nu'$ are different ones among them, the quantity

\[
\nu^2 - \nu'^2 = (\nu - \nu')(\nu + \nu')
\]
can never
%CONT>>
